% =============================================================================
%  ISM COMPANION READER  —  LaTeX template  (gold-standard faithful)
% =============================================================================
%  Reproduces Saurabh's hand-made companion PDF:
%    * fancyhdr running header ("ISM Companion" | "Session N · <Title>") + rule
%    * full-width DARK NAVY title banner (label / big title / subtitle / meta)
%    * numbered, sans, navy section headings with a rule beneath (titlesec)
%    * FIVE pill-tab tcolorbox callouts: intuition / everyday / worked /
%      watchout / keytake — each a rounded pill title tab on the top-left edge,
%      a leading glyph, a thin colour frame, and a very light tint fill.
%    * \housefig{path}{caption} for full-width matplotlib PNG figures.
%
%  BUILD:   pdflatex companion.tex && pdflatex companion.tex
%           (run twice so the table of contents / hyperref refs settle)
%
%  KEYLESS / OFFLINE: no network, no secrets. Real companions drop matplotlib
%  PNGs in via \housefig; this template's specimen uses an inline tikzpicture
%  so the FILE ALONE compiles to a sensible PDF with no PNG present yet.
%
%  HOST AGENT: replace every {{PLACEHOLDER}}. The placeholders that drive the
%  running header live in the three \newcommand lines just below \begin{document}
%  is set up in the preamble (see \CourseShort / \SessionLine / \LectureTitle).
%  One specimen section ("Bernoulli: the atom of randomness") is fully written,
%  using ALL FIVE callouts + a fully worked example, so the template renders
%  meaningfully. Duplicate its structure per slide-concept, then delete it.
% =============================================================================

\documentclass[a4paper,11pt]{article}

% ---- Encoding & fonts -------------------------------------------------------
\usepackage[T1]{fontenc}
\usepackage[utf8]{inputenc}
\usepackage{lmodern}
\usepackage{microtype}              % protrusion/kerning -> fewer overfull boxes

% ---- Geometry (~1in margins) ------------------------------------------------
\usepackage[a4paper,margin=1in,headheight=15pt,headsep=14pt]{geometry}

% ---- Math -------------------------------------------------------------------
\usepackage{amsmath,amssymb}

% ---- Graphics, tables, colour ----------------------------------------------
\usepackage{graphicx}
\usepackage{booktabs}
\usepackage[table,svgnames]{xcolor}

% ---- Callout boxes (breakable pill-tab tcolorboxes) -------------------------
\usepackage[most]{tcolorbox}
\tcbuselibrary{skins,breakable}

% ---- Tikz (ONLY for the specimen's inline PMF; real figs use \housefig) -----
\usepackage{tikz}

% ---- Callout glyphs (robust + cross-platform) -------------------------------
% pifont ALWAYS ships with TeX Live and gives crisp dingbats: ✗ (\ding{55}),
% ✓ (\ding{51}), the pointing hand (\ding{43}) and the pencil (\ding{46}).
% bbding's \Pointinghand / \Pencil look a touch nicer, but bbding is MISSING on
% some minimal installs (TinyTeX, and several Codex / Jules / CI sandboxes) — a
% hard "\usepackage{bbding}" then aborts the whole build. So load bbding ONLY if
% it is present and fall back to pifont otherwise. Use \glyphHand / \glyphPencil
% everywhere below so the callouts render identically either way.
\usepackage{pifont}
\IfFileExists{bbding.sty}{%
  \usepackage{bbding}%
  \newcommand{\glyphHand}{\Pointinghand}%
  \newcommand{\glyphPencil}{\Pencil}%
}{%
  \newcommand{\glyphHand}{\ding{43}}%
  \newcommand{\glyphPencil}{\ding{46}}%
}

% ---- Lists & spacing --------------------------------------------------------
\usepackage{enumitem}
\setlist{leftmargin=1.5em,topsep=4pt,itemsep=3pt}

% ---- Dedicated "Step 1. / Step 2. ..." list for WORKED EXAMPLES -------------
% WHY THIS EXISTS (a real bug it fixes): the bold "Step N." label is WIDE. With
% the small global leftmargin above, enumitem cannot fit the label in its box
% and overflows it to the LEFT — pushing "Step 1." OUTSIDE the callout's frame.
% This dedicated list reserves enough label width so every "Step N." sits
% cleanly INSIDE the box, with the step text aligned in a neat hanging column.
% ALWAYS use \begin{steps}...\end{steps} for worked-example steps.
% NEVER write a raw \begin{enumerate}[label=\textbf{Step \arabic*.}] — it leaks.
\newlist{steps}{enumerate}{1}
\setlist[steps]{label=\textbf{Step \arabic*.},
  leftmargin=4.4em, labelwidth=3.8em, labelsep=0.6em, labelindent=0pt,
  align=left, topsep=5pt, itemsep=6pt}
\usepackage{parskip}                % blank-line paragraph spacing, no indent

% ---- Headers / footers ------------------------------------------------------
\usepackage{fancyhdr}

% ---- Section styling --------------------------------------------------------
\usepackage{titlesec}

% ---- Links (hidden) + URL line-breaking so long URLs never overflow ---------
\usepackage[hidelinks]{hyperref}
\usepackage{xurl}
\hypersetup{breaklinks=true}

% =============================================================================
%  COLOURS — navy chrome + the FIVE taxonomy ramps (frame + light fill)
%  Frame hexes match the gold standard; fills approximate the tcolorbox
%  "blue!4 / orange!7 / green!5 / red!4 / violet!6" tints with explicit HTML so
%  the result is identical on every TeX distribution.
% =============================================================================
\definecolor{navy}{HTML}{21355E}        % banner + headings
\definecolor{navyrule}{HTML}{21355E}
\definecolor{inktext}{HTML}{1A202C}     % body ink
\definecolor{mutedtext}{HTML}{6B7280}   % header / meta grey

% Intuition — blue
\definecolor{intuFrame}{HTML}{2C5AA0}   \definecolor{intuFill}{HTML}{F4F7FC}  % ~ blue!4
% Everyday — amber
\definecolor{everyFrame}{HTML}{8A5A1E}  \definecolor{everyFill}{HTML}{FBF1E3}  % ~ orange!7
% Worked — green
\definecolor{workFrame}{HTML}{2E7D52}   \definecolor{workFill}{HTML}{F1F8F3}  % ~ green!5
% Watch out — red
\definecolor{watchFrame}{HTML}{B23A48}  \definecolor{watchFill}{HTML}{FBF1F2}  % ~ red!4
% Key takeaway — purple/violet
\definecolor{keyFrame}{HTML}{6A4C93}    \definecolor{keyFill}{HTML}{F4F0F9}    % ~ violet!6

\color{inktext}
\hypersetup{linkcolor=navy,urlcolor=navy,citecolor=navy}

% =============================================================================
%  RUNNING HEADER / FOOTER  (fancyhdr)  — every page
%  Left:  "{{COURSE_SHORT}}"      Right: "Session {{SESSION}} · {{LECTURE_TITLE}}"
%  A thin rule sits below; the footer is the centred page number.
%  Edit the three macros below (or let the host agent fill the {{...}}).
% =============================================================================
\newcommand{\CourseShort}{{{COURSE_SHORT}}}          % e.g. "ISM Companion"
\newcommand{\SessionNo}{{{SESSION}}}                 % e.g. "6"
\newcommand{\LectureTitle}{{{LECTURE_TITLE}}}        % e.g. "Probability Distributions"

\pagestyle{fancy}
\fancyhf{}
\fancyhead[L]{\footnotesize\color{mutedtext}\textsf{\CourseShort}}
\fancyhead[R]{\footnotesize\color{mutedtext}\textsf{Session \SessionNo\ \textperiodcentered\ \LectureTitle}}
\fancyfoot[L]{\footnotesize\color{mutedtext}\textsf{WILP, BITS Pilani}}
\fancyfoot[C]{\footnotesize\color{mutedtext}\thepage}
\renewcommand{\headrulewidth}{0.4pt}
\renewcommand{\footrulewidth}{0pt}
\renewcommand{\headrule}{\hbox to\headwidth{\color{navyrule!55}\leaders\hrule height \headrulewidth\hfill}}

% =============================================================================
%  SECTION HEADINGS (titlesec) — sans, bold, navy, NUMBERED, rule beneath
%  Produces e.g. "1  The big idea: ..." with a full-width navy hairline under it.
% =============================================================================
\titleformat{\section}
  {\sffamily\Large\bfseries\color{navy}}      % font
  {\thesection}                               % label (the number)
  {0.8em}                                      % gap label -> title
  {}                                           % before-code
  [{\vspace{2pt}\color{navy!70}\titlerule[0.8pt]}]   % rule beneath
\titleformat{\subsection}
  {\sffamily\large\bfseries\color{navy!92}}
  {\thesubsection}{0.7em}{}
\titlespacing*{\section}{0pt}{16pt}{8pt}
\titlespacing*{\subsection}{0pt}{12pt}{4pt}

% Each top-level TOPIC starts on a fresh page, so a heading never lands at the
% bottom of a page with its content spilling across the break. We skip the break
% before the FIRST section only — it continues under the title banner — using a
% one-shot flag, which is robust no matter when titlesec steps the section
% counter. Every section after the first clears to a new page.
\newif\ifmlkfirstsec \mlkfirstsectrue
\newcommand{\sectionbreak}{\ifmlkfirstsec\mlkfirstsecfalse\else\clearpage\fi}

% =============================================================================
%  PILL-TAB CALLOUTS  —  the FIVE taxonomy boxes
%  Shared look: enhanced + breakable, rounded corners, boxrule 0.6pt, very light
%  tint fill, thin colour frame, and a ROUNDED PILL boxed-title tab attached to
%  the top-left edge (xshift 6mm, yshift -3mm) in the frame colour with white
%  bold sans text and a leading glyph.
%
%  A single \housecallout style macro keeps all five visually identical; each
%  environment only swaps the frame colour, fill colour, and the pill text.
% =============================================================================
\tcbset{
  housecallout/.style 2 args={
    % #1 = frame/pill colour,  #2 = light fill colour
    enhanced, breakable,
    colback=#2, colframe=#1,
    boxrule=0.6pt, arc=2.4pt, outer arc=2.4pt,
    left=10pt, right=10pt, top=9pt, bottom=9pt,
    before skip=11pt, after skip=11pt,
    fontupper=\normalsize,
    coltitle=white,
    fonttitle=\bfseries\sffamily\small,
    attach boxed title to top left={xshift=6mm,yshift=-3mm},
    boxed title style={
      colback=#1, colframe=#1,
      rounded corners, arc=3pt,
      boxrule=0pt, left=6pt, right=6pt, top=2pt, bottom=2pt
    }
  }
}

% --- Intuition (blue, glyph ☞) ------------------------------------------------
\newtcolorbox{intuition}{%
  housecallout={intuFrame}{intuFill},
  title={\raisebox{-0.05ex}{\glyphHand}\,\ The intuition}}

% --- Everyday picture (amber, glyph ★) ---------------------------------------
\newtcolorbox{everyday}{%
  housecallout={everyFrame}{everyFill},
  title={\raisebox{-0.05ex}{$\bigstar$}\,\ Everyday picture}}

% --- Worked example (green, glyph ✎, takes a caption argument) -----------------
%  Usage: \begin{worked}{The caption}  ...steps...  \end{worked}
\newtcolorbox{worked}[1]{%
  housecallout={workFrame}{workFill},
  title={\raisebox{-0.05ex}{\glyphPencil}\,\ Worked example: #1}}

% --- Watch out (red, glyph ✗) -------------------------------------------------
\newtcolorbox{watchout}{%
  housecallout={watchFrame}{watchFill},
  title={\raisebox{-0.05ex}{\ding{55}}\,\ Watch out}}

% --- Key takeaway (purple, glyph ✓) ------------------------------------------
\newtcolorbox{keytake}{%
  housecallout={keyFrame}{keyFill},
  title={\raisebox{-0.05ex}{\ding{51}}\,\ Key takeaway}}

% =============================================================================
%  BITS BRAND  —  "WILP, BITS Pilani" wordmark + optional logo file
%  The institute brand is fixed on every kit. The white wordmark ALWAYS renders
%  (legible on the navy banner). If a logo image is present it is shown in a small
%  WHITE chip, so any logo colour stays readable on navy. To enable the logo, drop
%  a PNG at templates/assets/bits-logo.png (or output/<slug>/assets/bits-logo.png).
%  No file => wordmark only; nothing breaks. See templates/assets/README.md.
% =============================================================================
\newif\ifbitslogo \bitslogofalse
\IfFileExists{assets/bits-logo.png}%
  {\bitslogotrue \def\bitslogopath{assets/bits-logo.png}}%
  {\IfFileExists{../../templates/assets/bits-logo.png}%
     {\bitslogotrue \def\bitslogopath{../../templates/assets/bits-logo.png}}%
     {}}
% White wordmark — always shown.
\newcommand{\bitswordmark}{{\sffamily\bfseries\color{white}\fontsize{12.5}{15}\selectfont WILP, BITS Pilani}}
% Logo inside a white chip — only when a logo file was found (any colour stays legible).
\newcommand{\bitslogochip}{\ifbitslogo\colorbox{white}{\,\raisebox{-0.2ex}{\includegraphics[height=26pt]{\bitslogopath}}\,}\fi}

% =============================================================================
%  TITLE BANNER  —  full-width dark navy tcolorbox on page 1
%  \titlebanner{LABEL}{BIG TITLE}{SUBTITLE}{META LINE}
%    #1 small uppercase tracked label
%    #2 very large bold title
%    #3 one-line lighter subtitle
%    #4 small meta line (shown under a thin white rule)
%  Full text width via \linewidth; white text; breakable=false keeps it whole.
% =============================================================================
% \addfontfeatureslike is a tiny shim so the label field stays a plain
% pass-through even without a letterspacing package; defined as identity here so
% it is in scope before \titlebanner is ever expanded.
\providecommand{\addfontfeatureslike}[1]{#1}
\newcommand{\titlebanner}[4]{%
  \begin{tcolorbox}[
    enhanced, breakable=false,
    colback=navy, colframe=navy,
    boxrule=0pt, arc=3pt, outer arc=3pt,
    left=16pt, right=16pt, top=16pt, bottom=16pt,
    before skip=2pt, after skip=14pt,
    width=\linewidth ]
    % --- BITS brand row: wordmark left, optional logo chip right ---
    \noindent
    \begin{minipage}[c]{0.72\linewidth}\bitswordmark\end{minipage}%
    \hfill
    \begin{minipage}[c]{0.26\linewidth}\raggedleft\bitslogochip\end{minipage}\par
    \vspace{9pt}
    {\sffamily\footnotesize\bfseries\color{white!85}%
      \MakeUppercase{\addfontfeatureslike{#1}}}\par
    \vspace{6pt}
    {\sffamily\bfseries\fontsize{30}{34}\selectfont\color{white} #2\par}
    \vspace{5pt}
    {\sffamily\large\color{white!88} #3\par}
    \vspace{9pt}
    {\color{white!55}\rule{\linewidth}{0.4pt}}\par
    \vspace{6pt}
    {\sffamily\footnotesize\color{white!82} #4\par}
  \end{tcolorbox}%
}

% =============================================================================
%  \housefig{path}{caption}  —  full-width, centred matplotlib PNG figure
%  width=\linewidth guarantees the image never exceeds the text block (no
%  overfull \hbox). The plot's own baked-in bold title carries the headline; the
%  \caption gives the figure number + a short description.
% =============================================================================
\newcommand{\housefig}[2]{%
  \begin{figure}[htbp]
    \centering
    \includegraphics[width=\linewidth]{#1}%
    \caption{#2}
  \end{figure}%
}

% A small inline-vocabulary helper: bold-navy first-use term.
\newcommand{\vocab}[1]{\textbf{\color{navy}#1}}

% =============================================================================
%  OVERFLOW-SAFETY CHEAT-SHEET (read before filling)
% =============================================================================
%  * WIDE MATH: never one long line. Wrap with align / split / multline so it
%        breaks at the margin, e.g.
%            \begin{align} a &= b + c \\ &\quad + d \end{align}
%  * WIDE TABLE: \resizebox{\linewidth}{!}{\begin{tabular}...\end{tabular}}
%        (or booktabs columns sized with p{..} that sum under \linewidth).
%  * LONG URL: \url{...} (xurl lets it break anywhere).
%  * FIGURES: always \housefig (width=\linewidth) — never a raw oversized box.
%  * CALLOUTS: all breakable, so they flow cleanly to the next page; but keep a
%        single WORKED example coherent — don't let one example sprawl over a
%        page break mid-step if you can split the content into two boxes.
%  * WORKED-EXAMPLE STEPS: ALWAYS use \begin{steps}...\end{steps} (defined above),
%        never \begin{enumerate}[label=\textbf{Step \arabic*.}]. The steps list
%        reserves label width so "Step N." can never spill OUTSIDE the box frame.
%  * CONTENT INSIDE A BOX must fit the box: a wide table or long equation in a
%        callout still needs \resizebox / align — the box does not shrink it for
%        you. Anything wider than the box's text area overflows its frame.
% =============================================================================

\begin{document}

% ----------------------------------------------------------------------------
%  PAGE 1 — TITLE BANNER
%  AGENT:FILL — replace the four banner fields. Keep the meta line format:
%  "Session N (dates) · Read alongside the lecture slides · Every slide example
%   worked out in full".
% ----------------------------------------------------------------------------
\thispagestyle{fancy}

\titlebanner
  {Introduction to Statistical Methods \textperiodcentered\ Companion Reader}% LABEL
  {{{LECTURE_TITLE}}}                                                        % BIG TITLE
  {{{ONE_LINE_SUBTITLE}}}                                                    % SUBTITLE
  {Session \SessionNo\ ({{SESSION_DATES}}) \textperiodcentered\ Read alongside the lecture slides \textperiodcentered\ Every slide example worked out in full} % META

% ----------------------------------------------------------------------------
%  BITS credit line (institute brand — kept on every kit)
% ----------------------------------------------------------------------------
\smallskip
\noindent{\small\color{mutedtext}%
Prepared for the Work Integrated Learning Programmes (WILP), BITS Pilani.}
\smallskip

% ----------------------------------------------------------------------------
%  HOW TO USE THIS COMPANION  (intro paragraph right after the banner)
% ----------------------------------------------------------------------------
\noindent\textbf{How to use this companion.}
The slides give you the formulas; this reader gives you the \emph{why}. Every
slide concept is expanded into a full, plain-English explanation, and every
slide example is worked out in full, step by step. Colour-coded boxes guide you:
the \textcolor{intuFrame}{\textbf{blue}} ``intuition'' box states the core idea
in one breath; the \textcolor{everyFrame}{\textbf{amber}} ``everyday picture''
box ties it to real life before any math; the \textcolor{workFrame}{\textbf{green}}
``worked example'' box solves a problem with real numbers and no skipped steps;
the \textcolor{watchFrame}{\textbf{red}} ``watch out'' box flags the common
mistake; and the \textcolor{keyFrame}{\textbf{purple}} ``key takeaway'' box is
the one sentence to remember. Read it alongside the slides, in order.

\medskip

% ============================================================================
%  AGENT:FILL — OPENING SECTION(S)
%  Typically: an "agenda / what this session covers" section, then one
%  \section per slide-concept. Use \vocab{...} on each term's first use, wrap
%  display math in align/split, and add a \housefig per matplotlib plot.
%  Example skeleton (uncomment + fill, or duplicate the specimen below):
%
%  \section{What this session covers}
%  {{AGENDA_PROSE_AND_BULLETS}}
%
%  \section{ {{CONCEPT_TITLE}} }
%  {{EXPANDED_PLAIN_ENGLISH_EXPLANATION — define every symbol}}
%  \begin{intuition} {{...}} \end{intuition}
%  \begin{everyday}  {{relatable real-life analogy BEFORE the math}} \end{everyday}
%  {{build the formula step by step, every symbol named}}
%  \begin{worked}{ {{caption}} } {{every step, real numbers}} \end{worked}
%  \begin{watchout}  {{the classic mistake}} \end{watchout}
%  \begin{keytake}   {{the one sentence to remember}} \end{keytake}
%  \housefig{figs/{{plot}}.png}{ {{figure caption}} }
% ============================================================================

% ============================================================================
%  SPECIMEN SECTION (fully written so the template renders meaningfully).
%  Uses ALL FIVE callouts + one fully worked example. The figure here is an
%  INLINE tikzpicture (a tiny 2-bar PMF) so the TEMPLATE ALONE compiles with NO
%  PNG present. In real companions, replace it with \housefig + a matplotlib PNG.
%  Duplicate this \section per slide-concept, then DELETE the specimen.
% ============================================================================
\section{The big idea: Bernoulli, the atom of randomness}

Many real questions have just two answers. A coin lands heads or tails. An email
is spam or not spam. A patient tests positive or negative. The
\vocab{Bernoulli distribution} is the simplest model of exactly this: one trial
with two possible outcomes. We call one outcome a ``success'' (worth $1$) and the
other a ``failure'' (worth $0$). A single number, \vocab{$p$}, gives the
probability of success; the probability of failure is then $1-p$, because the two
must add up to $1$.

\begin{intuition}
A Bernoulli trial is one yes/no question with a fixed chance of ``yes.'' The whole
distribution is pinned down by a single number, $p$ --- the probability of a
``yes.'' Everything else (the average, the spread) follows from that one number.
\end{intuition}

\begin{everyday}
Picture a slightly bent coin that comes up heads $60\%$ of the time. You flip it
\emph{once}. That single flip is a Bernoulli trial with $p = 0.6$. ``Heads'' is
the success ($1$); ``tails'' is the failure ($0$). You are not asking ``how many
heads in ten flips'' yet --- just one flip, two possible outcomes, one chance $p$.
\end{everyday}

% ---- Build the formula step by step, naming every symbol --------------------
We write $X \sim \text{Bernoulli}(p)$ to mean ``$X$ is a Bernoulli random variable
with success probability $p$.'' The random variable $X$ takes the value $1$ on a
success and $0$ on a failure. Its \vocab{probability mass function (PMF)} --- the
rule that gives the probability of each value --- is

\begin{align}
  P(X = 1) &= p, \\
  P(X = 0) &= 1 - p .
\end{align}

These two lines can be folded into one tidy formula. For $x \in \{0,1\}$,

\begin{equation}
  P(X = x) \;=\; p^{x}\,(1-p)^{\,1-x}.
\end{equation}

Check it term by term. Put $x = 1$: the exponent on $p$ is $1$ and on $(1-p)$ is
$1-1 = 0$, so $P = p^{1}(1-p)^{0} = p$. Put $x = 0$: now $P = p^{0}(1-p)^{1} = 1-p$.
The single formula reproduces both lines exactly --- nothing skipped.

The \vocab{mean} (the long-run average value of $X$) and the \vocab{variance} (how
much $X$ spreads around that average) are

\begin{align}
  \mathbb{E}[X] &= 0\cdot(1-p) + 1\cdot p = p, \\
  \operatorname{Var}(X) &= \mathbb{E}[X^2] - \big(\mathbb{E}[X]\big)^2
     = p - p^2 = p(1-p).
\end{align}

Here $\mathbb{E}[X^2] = 0^2(1-p) + 1^2 p = p$, because squaring $0$ and $1$ leaves
them unchanged. So the variance is $p(1-p)$ --- biggest at $p = 0.5$ (maximum
uncertainty) and zero at $p = 0$ or $p = 1$ (no uncertainty at all).

% ---- Fully worked example ---------------------------------------------------
\begin{worked}{the bent coin with $p = 0.6$}
Let $X \sim \text{Bernoulli}(0.6)$, so a success (``heads'') has probability
$p = 0.6$.
\begin{steps}
  \item \textbf{Probabilities of each outcome.}
        $P(X=1) = p = 0.6$ and $P(X=0) = 1-p = 0.4$. They add to $1$, as they
        must.
  \item \textbf{Mean.} $\mathbb{E}[X] = p = 0.6$. On average, $0.6$ of a head per
        flip --- which is just the success probability, as expected.
  \item \textbf{Variance.} $\operatorname{Var}(X) = p(1-p) = 0.6 \times 0.4 = 0.24$.
  \item \textbf{Standard deviation.} $\sqrt{0.24} \approx 0.49$. This is the
        typical distance of a single outcome from the mean $0.6$ --- large,
        because one flip is very uncertain.
  \item \textbf{Sanity check with the folded PMF.} For $x = 1$:
        $p^{1}(1-p)^{0} = 0.6$. For $x = 0$: $p^{0}(1-p)^{1} = 0.4$. Same answers
        as Step~1, so the single formula is consistent.
\end{steps}
\end{worked}

% ---- Specimen figure: INLINE tikz PMF (no PNG needed for the template) ------
% In a REAL companion this becomes a matplotlib PNG via, e.g.:
%     \housefig{figs/bernoulli_pmf.png}{Bernoulli($p{=}0.6$): one shot, two outcomes.}
% The inline tikz below keeps THIS template compilable with no PNG present.
\begin{figure}[htbp]
  \centering
  \begin{tikzpicture}[x=2.4cm,y=4cm]
    % axes
    \draw[->,gray] (-0.15,0) -- (1.5,0) node[right,black,font=\small]{$x$};
    \draw[->,gray] (0,-0.02) -- (0,1.05) node[above,black,font=\small]{$P(X=x)$};
    % gridline at 1.0
    \draw[gray!30] (0,1) -- (1.4,1) node[right,gray,font=\scriptsize]{$1.0$};
    % bar for x = 0  (height 0.4, failure, red ramp)
    \fill[watchFrame!75] (0.0,0) rectangle (0.28,0.4);
    \node[below,font=\small] at (0.14,0) {$0$};
    \node[above,font=\small] at (0.14,0.4) {$0.4$};
    % bar for x = 1  (height 0.6, success, green ramp)
    \fill[workFrame!80] (1.0,0) rectangle (1.28,0.6);
    \node[below,font=\small] at (1.14,0) {$1$};
    \node[above,font=\small] at (1.14,0.6) {$0.6$};
    % baked-in bold title, house style
    \node[font=\sffamily\bfseries\small,navy] at (0.7,1.18)
      {Bernoulli($p=0.6$): one shot, two outcomes};
  \end{tikzpicture}
  \caption{The Bernoulli($p=0.6$) PMF: failure ($x=0$) has mass $0.4$, success
  ($x=1$) has mass $0.6$. Inline placeholder --- real companions use a
  matplotlib PNG through \texttt{\textbackslash housefig}.}
\end{figure}

\begin{watchout}
$p$ is the probability of a \emph{single} success, not a count of successes. A
Bernoulli trial is \textbf{one} flip, so $X$ is only ever $0$ or $1$ --- never $2$
or $7$. As soon as you flip the coin many times and count the heads, you have left
Bernoulli and entered the \vocab{Binomial} distribution (the sum of many
independent Bernoulli trials). Don't mix the two up.
\end{watchout}

\begin{keytake}
A Bernoulli trial is one yes/no event with success probability $p$: its PMF is
$p^{x}(1-p)^{1-x}$, its mean is $p$, and its variance is $p(1-p)$. It is the
single building block from which the Binomial, and much of classification in
machine learning, is assembled.
\end{keytake}

% ---- Optional ML/AI connection (the teacher asks for one where natural) -----
\noindent\textbf{ML/AI connection.}
Every binary classifier --- spam vs.\ not-spam, fraud vs.\ legitimate --- models
its label as a Bernoulli variable. The model outputs $\hat p$, the predicted
probability of the positive class, and the training loss
(\vocab{binary cross-entropy}) is exactly the negative log-likelihood of a
Bernoulli: $-\big[y\log\hat p + (1-y)\log(1-\hat p)\big]$. So this ``atom of
randomness'' sits at the heart of logistic regression and the final layer of most
neural-network classifiers.

% ============================================================================
%  AGENT:FILL — REMAINING CONCEPT SECTIONS
%  Duplicate the specimen \section structure for each remaining slide-concept:
%  expand the slide's bullets/definition/formula into plain English, build the
%  math step by step, add the relevant callouts (intuition / everyday / worked /
%  watchout / keytake), an ML/AI connection where natural, and a \housefig per
%  matplotlib plot. Then DELETE the specimen above. Examples to template next:
%
%  \section{ {{CONCEPT_2_TITLE}} }   % e.g. "Binomial: counting the successes"
%  \section{ {{CONCEPT_3_TITLE}} }   % e.g. "Poisson: rare events over time"
%  ... one \section per slide-concept, in slide order ...
% ============================================================================

% ============================================================================
%  AGENT:FILL — CLOSE WITH A GLOSSARY + SYMBOL CHEAT-SHEET (companion_style §8)
%  A scannable revision page: every term in one plain line, every symbol with its
%  meaning. List only what the lecture used, in teaching order. Pattern:
%
%  \section*{Glossary \& symbols}
%  \addcontentsline{toc}{section}{Glossary \& symbols}
%  \noindent\textbf{Key terms.}
%  \begin{description}[leftmargin=8.5em,style=nextline,font=\normalfont\bfseries\color{navy}]
%    \item[term] one plain-English line of meaning.
%  \end{description}
%  \medskip
%  \noindent\textbf{Symbols at a glance.}
%  \begin{center}\small
%  \begin{tabular}{@{}ll@{}}
%  \toprule \textbf{Symbol} & \textbf{Meaning} \\ \midrule
%  $\sigma(z)$ & the sigmoid; squashes any number into (0,1) \\
%  \bottomrule
%  \end{tabular}
%  \end{center}
% ============================================================================

\end{document}
